%% kapa: exam=1 round=31 year=2020 period=1 form=A subject=economics subject_ko=경제학원론 seq=18 number=58
%% source: Q-Net 감정평가사 기출문제 (cst00309) artlSeq=5207219 file=extracted/1008/5207219/A형/A형/2020년도 제31회 감정평가사 1차 시험_1교시_A형.pdf
%% conversion: kordoc 4.24.51 -> tools/gen1.py (auto); review=proofread-sonnet; flags=none
\begin{question}{58}
두 재화 $X$ 와 $Y$ 를 소비하는 소비자 甲과 乙이 존재하는 순수교환경제를 가정한다. 두 소비자의 효용함수는 $U(x,y)=xy$ 로 동일하고, 甲의 초기부존은 ( $x=10,y=5$), 乙의 초기부존은 ( $x=5,y=10$)일 때, 옳은 것을 모두 고른 것은? (단, $x$는 $X$재 소비량, $y$는 $Y$재 소비량)
\begin{framedtext}
ㄱ. 초기부존에서 甲의 한계대체율은 0.5, 乙의 한계대체율은 2이다. \par
ㄴ. 초기부존에서 甲의 $X$재 1단위와 乙의 $Y$재 2단위가 교환될 때 파레토 개선이 이루어진다. \par
ㄷ. 일반균형은 $X$재 가격이 1일 때, $Y$재 가격은 1이다. \par
ㄹ. 일반균형에서 甲은 $X$재보다 $Y$재를 더 많이 소비한다. \par
\end{framedtext}
\begin{choices}
  \item ㄱ, ㄴ
  \item ㄱ, ㄷ
  \item ㄴ, ㄷ
  \item ㄴ, ㄹ
  \item ㄷ, ㄹ
\end{choices}
\end{question}
\answer{2}
